% ============================================================
% GLOSSARIO - BugBoard26
% Includere con % ============================================================
% GLOSSARIO - BugBoard26
% Includere con % ============================================================
% GLOSSARIO - BugBoard26
% Includere con % ============================================================
% GLOSSARIO - BugBoard26
% Includere con \input{glossario}
% ============================================================

\section{Glossario}

Di seguito sono riportati i termini tecnici e di dominio utilizzati nel presente documento e nel sistema BugBoard26.

\begin{description}[style=nextline, leftmargin=3cm, labelwidth=2.8cm]

    \item[Admin]
    Utente con privilegi amministrativi completi. Può gestire utenti, accedere a tutti i bug del sistema, generare report e modificare configurazioni.

    \item[Allegato]
    File immagine (JPEG, PNG, GIF, WebP, dimensione massima 5\,MB) caricato in associazione a un bug per documentarne visivamente il problema o la soluzione.

    \item[Archiviazione]
    Operazione che sposta un bug in stato \textit{Archived}, rimuovendolo dalla lista attiva senza eliminarlo definitivamente dal sistema.

    \item[Assegnatario]
    Utente registrato responsabile della risoluzione di un determinato bug.

    \item[Bug]
    Segnalazione di un comportamento anomalo, errato o inatteso di un sistema software. Nel contesto di BugBoard26 il termine indica genericamente qualsiasi segnalazione inserita nel sistema, indipendentemente dal tipo (Bug, Feature, Enhancement).

    \item[BugBoard26]
    Il sistema software oggetto di questo documento. Applicazione web per il tracciamento e la gestione collaborativa delle segnalazioni di bug.

    \item[Caso d'uso]
    Descrizione di un'interazione tra un attore e il sistema per raggiungere un obiettivo specifico. Formalizzato secondo il metodo di Alistair Cockburn.

    \item[Commento]
    Testo inserito da un utente in associazione a un bug per fornire aggiornamenti, chiarimenti o informazioni aggiuntive.

    \item[Docker]
    Piattaforma di containerizzazione utilizzata per impacchettare e distribuire i componenti dell'applicazione (frontend e backend) in modo indipendente dall'ambiente host.

    \item[Enhancement]
    Tipo di segnalazione che indica una richiesta di miglioramento di una funzionalità già esistente. Distinto da un Bug (errore) e da una Feature (nuova funzionalità).

    \item[Feature]
    Tipo di segnalazione che indica la richiesta di implementazione di una nuova funzionalità nel sistema.

    \item[History (Storico)]
    Registro cronologico di tutte le modifiche apportate a un bug nel tempo: cambi di stato, assegnazioni, commenti, aggiunta di allegati.

    \item[JWT (JSON Web Token)]
    Standard aperto (RFC 7519) utilizzato per autenticare in modo sicuro le richieste HTTP tra frontend e backend. Il token viene generato al login e incluso in ogni richiesta successiva.

    \item[Label]
    Etichetta testuale opzionale associabile a un bug per categorizzarlo ulteriormente (es.\ ``frontend'', ``performance'', ``regression'').

    \item[Priorità]
    Livello di urgenza di un bug. I valori previsti sono: \textit{Low}, \textit{Medium}, \textit{High}, \textit{Critical}.

    \item[ReadOnly (Utente)]
    Ruolo con accesso in sola lettura. L'utente ReadOnly può visualizzare bug, commenti e allegati, ma non può creare, modificare o commentare.

    \item[Report Mensile]
    Documento generato dal sistema contenente statistiche aggregate sull'attività del periodo selezionato: numero di bug creati, risolti, tempo medio di risoluzione, distribuzione per priorità e tipo.

    \item[Ruolo]
    Insieme di permessi associato a un utente. I ruoli previsti sono: \textit{USER} (utente standard), \textit{READONLY} (sola lettura), \textit{ADMIN} (amministratore).

    \item[Stato (Status)]
    Fase del ciclo di vita di un bug. Gli stati previsti sono:
    \begin{itemize}
        \item \textit{Todo}: bug appena creato, non ancora preso in carico
        \item \textit{In Progress}: bug in fase di risoluzione
        \item \textit{Resolved}: bug risolto
        \item \textit{Archived}: bug archiviato
    \end{itemize}

    \item[Tipo (Type)]
    Categoria della segnalazione. I valori previsti sono: \textit{Bug}, \textit{Feature}, \textit{Enhancement}.

    \item[USER (Utente registrato)]
    Ruolo standard. L'utente può creare bug, commentare, caricare allegati, assegnare bug (se Admin) e gestire il proprio profilo.

\end{description}

% ============================================================

\section{Glossario}

Di seguito sono riportati i termini tecnici e di dominio utilizzati nel presente documento e nel sistema BugBoard26.

\begin{description}[style=nextline, leftmargin=3cm, labelwidth=2.8cm]

    \item[Admin]
    Utente con privilegi amministrativi completi. Può gestire utenti, accedere a tutti i bug del sistema, generare report e modificare configurazioni.

    \item[Allegato]
    File immagine (JPEG, PNG, GIF, WebP, dimensione massima 5\,MB) caricato in associazione a un bug per documentarne visivamente il problema o la soluzione.

    \item[Archiviazione]
    Operazione che sposta un bug in stato \textit{Archived}, rimuovendolo dalla lista attiva senza eliminarlo definitivamente dal sistema.

    \item[Assegnatario]
    Utente registrato responsabile della risoluzione di un determinato bug.

    \item[Bug]
    Segnalazione di un comportamento anomalo, errato o inatteso di un sistema software. Nel contesto di BugBoard26 il termine indica genericamente qualsiasi segnalazione inserita nel sistema, indipendentemente dal tipo (Bug, Feature, Enhancement).

    \item[BugBoard26]
    Il sistema software oggetto di questo documento. Applicazione web per il tracciamento e la gestione collaborativa delle segnalazioni di bug.

    \item[Caso d'uso]
    Descrizione di un'interazione tra un attore e il sistema per raggiungere un obiettivo specifico. Formalizzato secondo il metodo di Alistair Cockburn.

    \item[Commento]
    Testo inserito da un utente in associazione a un bug per fornire aggiornamenti, chiarimenti o informazioni aggiuntive.

    \item[Docker]
    Piattaforma di containerizzazione utilizzata per impacchettare e distribuire i componenti dell'applicazione (frontend e backend) in modo indipendente dall'ambiente host.

    \item[Enhancement]
    Tipo di segnalazione che indica una richiesta di miglioramento di una funzionalità già esistente. Distinto da un Bug (errore) e da una Feature (nuova funzionalità).

    \item[Feature]
    Tipo di segnalazione che indica la richiesta di implementazione di una nuova funzionalità nel sistema.

    \item[History (Storico)]
    Registro cronologico di tutte le modifiche apportate a un bug nel tempo: cambi di stato, assegnazioni, commenti, aggiunta di allegati.

    \item[JWT (JSON Web Token)]
    Standard aperto (RFC 7519) utilizzato per autenticare in modo sicuro le richieste HTTP tra frontend e backend. Il token viene generato al login e incluso in ogni richiesta successiva.

    \item[Label]
    Etichetta testuale opzionale associabile a un bug per categorizzarlo ulteriormente (es.\ ``frontend'', ``performance'', ``regression'').

    \item[Priorità]
    Livello di urgenza di un bug. I valori previsti sono: \textit{Low}, \textit{Medium}, \textit{High}, \textit{Critical}.

    \item[ReadOnly (Utente)]
    Ruolo con accesso in sola lettura. L'utente ReadOnly può visualizzare bug, commenti e allegati, ma non può creare, modificare o commentare.

    \item[Report Mensile]
    Documento generato dal sistema contenente statistiche aggregate sull'attività del periodo selezionato: numero di bug creati, risolti, tempo medio di risoluzione, distribuzione per priorità e tipo.

    \item[Ruolo]
    Insieme di permessi associato a un utente. I ruoli previsti sono: \textit{USER} (utente standard), \textit{READONLY} (sola lettura), \textit{ADMIN} (amministratore).

    \item[Stato (Status)]
    Fase del ciclo di vita di un bug. Gli stati previsti sono:
    \begin{itemize}
        \item \textit{Todo}: bug appena creato, non ancora preso in carico
        \item \textit{In Progress}: bug in fase di risoluzione
        \item \textit{Resolved}: bug risolto
        \item \textit{Archived}: bug archiviato
    \end{itemize}

    \item[Tipo (Type)]
    Categoria della segnalazione. I valori previsti sono: \textit{Bug}, \textit{Feature}, \textit{Enhancement}.

    \item[USER (Utente registrato)]
    Ruolo standard. L'utente può creare bug, commentare, caricare allegati, assegnare bug (se Admin) e gestire il proprio profilo.

\end{description}

% ============================================================

\section{Glossario}

Di seguito sono riportati i termini tecnici e di dominio utilizzati nel presente documento e nel sistema BugBoard26.

\begin{description}[style=nextline, leftmargin=3cm, labelwidth=2.8cm]

    \item[Admin]
    Utente con privilegi amministrativi completi. Può gestire utenti, accedere a tutti i bug del sistema, generare report e modificare configurazioni.

    \item[Allegato]
    File immagine (JPEG, PNG, GIF, WebP, dimensione massima 5\,MB) caricato in associazione a un bug per documentarne visivamente il problema o la soluzione.

    \item[Archiviazione]
    Operazione che sposta un bug in stato \textit{Archived}, rimuovendolo dalla lista attiva senza eliminarlo definitivamente dal sistema.

    \item[Assegnatario]
    Utente registrato responsabile della risoluzione di un determinato bug.

    \item[Bug]
    Segnalazione di un comportamento anomalo, errato o inatteso di un sistema software. Nel contesto di BugBoard26 il termine indica genericamente qualsiasi segnalazione inserita nel sistema, indipendentemente dal tipo (Bug, Feature, Enhancement).

    \item[BugBoard26]
    Il sistema software oggetto di questo documento. Applicazione web per il tracciamento e la gestione collaborativa delle segnalazioni di bug.

    \item[Caso d'uso]
    Descrizione di un'interazione tra un attore e il sistema per raggiungere un obiettivo specifico. Formalizzato secondo il metodo di Alistair Cockburn.

    \item[Commento]
    Testo inserito da un utente in associazione a un bug per fornire aggiornamenti, chiarimenti o informazioni aggiuntive.

    \item[Docker]
    Piattaforma di containerizzazione utilizzata per impacchettare e distribuire i componenti dell'applicazione (frontend e backend) in modo indipendente dall'ambiente host.

    \item[Enhancement]
    Tipo di segnalazione che indica una richiesta di miglioramento di una funzionalità già esistente. Distinto da un Bug (errore) e da una Feature (nuova funzionalità).

    \item[Feature]
    Tipo di segnalazione che indica la richiesta di implementazione di una nuova funzionalità nel sistema.

    \item[History (Storico)]
    Registro cronologico di tutte le modifiche apportate a un bug nel tempo: cambi di stato, assegnazioni, commenti, aggiunta di allegati.

    \item[JWT (JSON Web Token)]
    Standard aperto (RFC 7519) utilizzato per autenticare in modo sicuro le richieste HTTP tra frontend e backend. Il token viene generato al login e incluso in ogni richiesta successiva.

    \item[Label]
    Etichetta testuale opzionale associabile a un bug per categorizzarlo ulteriormente (es.\ ``frontend'', ``performance'', ``regression'').

    \item[Priorità]
    Livello di urgenza di un bug. I valori previsti sono: \textit{Low}, \textit{Medium}, \textit{High}, \textit{Critical}.

    \item[ReadOnly (Utente)]
    Ruolo con accesso in sola lettura. L'utente ReadOnly può visualizzare bug, commenti e allegati, ma non può creare, modificare o commentare.

    \item[Report Mensile]
    Documento generato dal sistema contenente statistiche aggregate sull'attività del periodo selezionato: numero di bug creati, risolti, tempo medio di risoluzione, distribuzione per priorità e tipo.

    \item[Ruolo]
    Insieme di permessi associato a un utente. I ruoli previsti sono: \textit{USER} (utente standard), \textit{READONLY} (sola lettura), \textit{ADMIN} (amministratore).

    \item[Stato (Status)]
    Fase del ciclo di vita di un bug. Gli stati previsti sono:
    \begin{itemize}
        \item \textit{Todo}: bug appena creato, non ancora preso in carico
        \item \textit{In Progress}: bug in fase di risoluzione
        \item \textit{Resolved}: bug risolto
        \item \textit{Archived}: bug archiviato
    \end{itemize}

    \item[Tipo (Type)]
    Categoria della segnalazione. I valori previsti sono: \textit{Bug}, \textit{Feature}, \textit{Enhancement}.

    \item[USER (Utente registrato)]
    Ruolo standard. L'utente può creare bug, commentare, caricare allegati, assegnare bug (se Admin) e gestire il proprio profilo.

\end{description}

% ============================================================

\section{Glossario}

Di seguito sono riportati i termini tecnici e di dominio utilizzati nel presente documento e nel sistema BugBoard26.

\begin{description}[style=nextline, leftmargin=3cm, labelwidth=2.8cm]

    \item[Admin]
    Utente con privilegi amministrativi completi. Può gestire utenti, accedere a tutti i bug del sistema, generare report e modificare configurazioni.

    \item[Allegato]
    File immagine (JPEG, PNG, GIF, WebP, dimensione massima 5\,MB) caricato in associazione a un bug per documentarne visivamente il problema o la soluzione.

    \item[Archiviazione]
    Operazione che sposta un bug in stato \textit{Archived}, rimuovendolo dalla lista attiva senza eliminarlo definitivamente dal sistema.

    \item[Assegnatario]
    Utente registrato responsabile della risoluzione di un determinato bug.

    \item[Bug]
    Segnalazione di un comportamento anomalo, errato o inatteso di un sistema software. Nel contesto di BugBoard26 il termine indica genericamente qualsiasi segnalazione inserita nel sistema, indipendentemente dal tipo (Bug, Feature, Enhancement).

    \item[BugBoard26]
    Il sistema software oggetto di questo documento. Applicazione web per il tracciamento e la gestione collaborativa delle segnalazioni di bug.

    \item[Caso d'uso]
    Descrizione di un'interazione tra un attore e il sistema per raggiungere un obiettivo specifico. Formalizzato secondo il metodo di Alistair Cockburn.

    \item[Commento]
    Testo inserito da un utente in associazione a un bug per fornire aggiornamenti, chiarimenti o informazioni aggiuntive.

    \item[Docker]
    Piattaforma di containerizzazione utilizzata per impacchettare e distribuire i componenti dell'applicazione (frontend e backend) in modo indipendente dall'ambiente host.

    \item[Enhancement]
    Tipo di segnalazione che indica una richiesta di miglioramento di una funzionalità già esistente. Distinto da un Bug (errore) e da una Feature (nuova funzionalità).

    \item[Feature]
    Tipo di segnalazione che indica la richiesta di implementazione di una nuova funzionalità nel sistema.

    \item[History (Storico)]
    Registro cronologico di tutte le modifiche apportate a un bug nel tempo: cambi di stato, assegnazioni, commenti, aggiunta di allegati.

    \item[JWT (JSON Web Token)]
    Standard aperto (RFC 7519) utilizzato per autenticare in modo sicuro le richieste HTTP tra frontend e backend. Il token viene generato al login e incluso in ogni richiesta successiva.

    \item[Label]
    Etichetta testuale opzionale associabile a un bug per categorizzarlo ulteriormente (es.\ ``frontend'', ``performance'', ``regression'').

    \item[Priorità]
    Livello di urgenza di un bug. I valori previsti sono: \textit{Low}, \textit{Medium}, \textit{High}, \textit{Critical}.

    \item[ReadOnly (Utente)]
    Ruolo con accesso in sola lettura. L'utente ReadOnly può visualizzare bug, commenti e allegati, ma non può creare, modificare o commentare.

    \item[Report Mensile]
    Documento generato dal sistema contenente statistiche aggregate sull'attività del periodo selezionato: numero di bug creati, risolti, tempo medio di risoluzione, distribuzione per priorità e tipo.

    \item[Ruolo]
    Insieme di permessi associato a un utente. I ruoli previsti sono: \textit{USER} (utente standard), \textit{READONLY} (sola lettura), \textit{ADMIN} (amministratore).

    \item[Stato (Status)]
    Fase del ciclo di vita di un bug. Gli stati previsti sono:
    \begin{itemize}
        \item \textit{Todo}: bug appena creato, non ancora preso in carico
        \item \textit{In Progress}: bug in fase di risoluzione
        \item \textit{Resolved}: bug risolto
        \item \textit{Archived}: bug archiviato
    \end{itemize}

    \item[Tipo (Type)]
    Categoria della segnalazione. I valori previsti sono: \textit{Bug}, \textit{Feature}, \textit{Enhancement}.

    \item[USER (Utente registrato)]
    Ruolo standard. L'utente può creare bug, commentare, caricare allegati, assegnare bug (se Admin) e gestire il proprio profilo.

\end{description}
